\section{Euclidean shapes and Deligne--Mostow periods}
\label{sec:thurstonDM}

\subsection{The fixed-curvature leaf at the Euclidean end}

Fix $n\ge4$ and real numbers
\begin{equation}
 0<\mu_i<1,\qquad \sum_{i=1}^n\mu_i=2.
 \label{eq:flatweights}
\end{equation}
The Euclidean cone opening at the $i$th marked point is
$2\pi(1-\mu_i)$, so $2\pi\mu_i$ is its curvature deficit. A conformal
sphere $X$ with these marked points carries a unique such flat metric up
to multiplication by a positive constant. In a coordinate with all marks
finite it has the explicit form
\begin{equation}
 g_{\mathrm E,\mu}=C\prod_i|w-p_i|^{-2\mu_i}|\dd w|^2.
 \label{eq:flatconemetric}
\end{equation}
The sum condition makes infinity a regular point. Near $p_i$, radial
integration gives a cone of opening $2\pi(1-\mu_i)$. Thus the space of
oriented, labeled Euclidean cone spheres with these deficits, up to
similarity, has underlying complex manifold $\M_{0,n}$.

Alexandrov realization identifies these metrics with convex Euclidean
polyhedral surfaces, allowing doubled planar polygons. This is the
fixed-deficit space on which Thurston constructs his complex hyperbolic
metric~\cite{ThurstonShapes,Fillastre}. The qualification ``fixed
deficits'' is essential: all Euclidean polyhedra with $n$ vertices, up
to similarity, form a space of real dimension $3n-7$, whereas the
fixed-deficit leaf has real dimension $2n-6$. The full space of compact
hyperbolic $n$-vertex polyhedra has dimension $3n-6$; its fixed-cone-angle
leaves again have dimension $2n-6$. The ideal locus has dimension $2n-6$.
These dimension counts apply on the generic strata; changes of
combinatorics and flattened realizations require the usual cell structure.


% BEGIN INLINE FIGURE: 02_polyhedral_regimes.tex
\begin{figure}[htbp]
\centering
\begin{tikzpicture}[font=\small,line cap=round,line join=round]
  \begin{scope}[shift={(-3.35,2.15)}]
    \node[align=center,font=\bfseries\small] at (0,1.95) {Euclidean tangent geometry};
    \draw[mpblue!15,line width=.5pt] (-1.7,-1.1) grid[step=.55] (1.7,1.15);
    \filldraw[fill=mpblue!8,draw=mpblue,line width=1pt]
      (-1.1,-.65)--(1.1,-.65)--(.05,1.1)--cycle;
    \foreach \p in {(-1.1,-.65),(1.1,-.65),(.05,1.1)}\fill[mpblue] \p circle (2pt);
    \node[align=center,text width=4.5cm,text=mpgray] at (0,-1.65)
      {Shrink, then rescale\\about an interior point};
  \end{scope}
  \begin{scope}[shift={(3.35,2.15)}]
    \node[font=\bfseries\small] at (0,1.95) {Compact vertices};
    \draw[mpgray!65,line width=.6pt] (0,0) circle (1.35);
    \filldraw[fill=mpblue!8,draw=mpblue,line width=1pt]
      (90:.9)--(210:.9)--(330:.9)--cycle;
    \foreach \a in {90,210,330}\fill[mpblue] (\a:.9) circle (2pt);
    \node[text=mpgray] at (0,-1.65) {Vertices inside the Klein ball};
  \end{scope}
  \begin{scope}[shift={(-3.35,-2.35)}]
    \node[font=\bfseries\small] at (0,1.95) {Ideal vertices};
    \draw[mpgray!65,line width=.6pt] (0,0) circle (1.35);
    \filldraw[fill=mpteal!8,draw=mpteal,line width=1pt]
      (90:1.35)--(210:1.35)--(330:1.35)--cycle;
    \foreach \a in {90,210,330}\filldraw[fill=white,draw=mpteal,line width=1pt] (\a:1.35) circle (2.4pt);
    \node[text=mpgray] at (0,-1.65) {Vertices on the sphere at infinity};
  \end{scope}
  \begin{scope}[shift={(3.35,-2.35)}]
    \node[font=\bfseries\small] at (0,1.95) {Hyperideal vertices};
    \draw[mpgray!65,line width=.6pt] (0,0) circle (1.35);
    \draw[mpgray!60,dashed] (90:1.55)--(210:1.55)--(330:1.55)--cycle;
    \fill[mpteal!8]
      (-.21604075,1.17580645)--(-1.12629863,-.40080645)--(-.91025788,-.775)
      --(.91025788,-.775)--(1.12629863,-.40080645)--(.21604075,1.17580645)--cycle;
    \draw[mpblue,line width=1pt]
      (-.21604075,1.17580645)--(-1.12629863,-.40080645)
      (-.91025788,-.775)--(.91025788,-.775)
      (1.12629863,-.40080645)--(.21604075,1.17580645);
    \draw[mpochre,line width=1.4pt]
      (-.21604075,1.17580645)--(.21604075,1.17580645)
      (-1.12629863,-.40080645)--(-.91025788,-.775)
      (.91025788,-.775)--(1.12629863,-.40080645);
    \foreach \a in {90,210,330}\fill[mpgray] (\a:1.55) circle (1.6pt);
    \node[text=mpgray,align=center] at (0,-1.65)
      {Exterior vertices and\\polar truncation planes};
  \end{scope}
\end{tikzpicture}
\caption[Euclidean, compact, ideal, and hyperideal regimes]{Local pictures for the geometric regimes. The three Klein-model
drawings are planar sections illustrating the location of vertices;
the ochre segments are polar truncations in the hyperideal case.
The Euclidean picture represents a tangent-scale limit after shrinking
and rescaling. These drawings do not assert a single global deformation
through all four regimes with arbitrary prescribed data.}
\label{fig:mono-polyhedral-regimes}
\end{figure}
% END INLINE FIGURE

\subsection{The same periods describe hypergeometric functions and developed edges}

Define the multivalued differential
\begin{equation}
 \eta_p(w)=\prod_i(w-p_i)^{-\mu_i}\,\dd w.
 \label{eq:DMdifferential}
\end{equation}
Its branches differ by unit complex numbers. It therefore defines a
twisted cohomology class for a rank-one unitary local system on the
punctured sphere. The relevant twisted cohomology space has complex
dimension $n-2$, and its Hermitian intersection form $H$ has signature
$(1,n-3)$, with the convention that the class of $\eta_p$ has positive
norm. In a local basis of twisted cycles or regularized relative arcs,
the coordinates
\begin{equation}
 Z_j(p)=\int_{\gamma_j}\eta_p
 \label{eq:DMperiods}
\end{equation}
are the hypergeometric periods of Deligne--Mostow~\cite{DeligneMostow}.
Projectivizing their positive line gives a local map to
$\mathbb{CH}^{n-3}$. Analytic continuation produces projective unitary
monodromy.

Thurston's Euclidean construction uses exactly the same coordinates.
The local developing map is
\[
 F(w)=\int^w\eta_p.
\]
Hence a developed edge vector is an integral of $\eta_p$ along the
corresponding arc. Cutting along a tree produces a Euclidean polygon
whose edge vectors are periods, and whose area is the Hermitian form
in these periods. Normalize $H$ by this area pairing. Then, more
intrinsically,
\begin{equation}
 \mathcal A(p)=\frac{i}{2}\int_{\mathbb P^1}\eta_p\wedge\overline{\eta_p}
             =H(Z(p),Z(p)).
 \label{eq:DMarea}
\end{equation}
Multiplication of all periods by a nonzero complex number changes only
the scale and rotation of the developed surface. Quotienting by this
action is precisely the projectivization in the Deligne--Mostow
construction.

Thus the connection is an equality of geometric structures: the
hypergeometric periods are the developed edge vectors, and the
Hermitian intersection pairing is the area pairing. Kojima proves that
the two constructions, with labels retained, are canonically isometric
\cite[Lemmas 4.1--4.2 and Theorem 4.3]{Kojima}. The scalar used for
the Hermitian intersection pairing varies between sources; fixing it
by~\eqref{eq:DMarea} removes that ambiguity, and a constant scalar
does not change the projective metric. Throughout this monograph the complex hyperbolic
metric has \emph{holomorphic sectional curvature $-1$}. Its K\"ahler
form, in the preceding local period coordinates, is
\begin{equation}
 \omega_{\mathrm{CH}}=-2i\,\partial\bar\partial\log\mathcal A.
 \label{eq:CHpotential}
\end{equation}
A holomorphic change of local scale for $\eta_p$ does not affect this
form. On a complex line in the ball, $\mathcal A=1-|z|^2$ gives
$\dd s^2=4|\dd z|^2/(1-|z|^2)^2$, which checks the curvature
normalization directly. A curvature-$-4$ convention would divide the form by $4$ and
the volume in complex dimension $d$ by $4^d$.

If the weights are rational, choose their least common denominator
$q$, so $\mu_i=m_i/q$. The compact normalization of the connected cyclic
cover
\[
 y^q=\prod_i(w-p_i)^{m_i}
\]
makes $\eta_p=\dd w/y$ a single-valued holomorphic differential. Under
the deck transformation $y\mapsto e^{2\pi i/q}y$ it has eigenvalue
$e^{-2\pi i/q}$. The corresponding eigenspace in the ordinary
cohomology of this cover realizes the twisted cohomology. This gives a
further useful interpretation in terms of periods of algebraic curves;
it does not require the monodromy group to be a lattice.

\subsection{Collision angles and the lattice conditions}

Suppose two marks of total weight $\mu_i+\mu_j<1$ collide, with local
separation $\delta$, while the other marks remain distinct and away
from the collision. After fixing a branch, rescaling the integration variable in their
vanishing period gives
\begin{equation}
 Z_{ij}(\delta)=C\delta^{\,1-\mu_i-\mu_j}(1+O(\delta)),\qquad C\ne0.
 \label{eq:vanishingperiod}
\end{equation}
The transverse complex hyperbolic cone angle is consequently
$2\pi(1-\mu_i-\mu_j)$. Deligne--Mostow's condition INT is
\begin{equation}
 (1-\mu_i-\mu_j)^{-1}\in\ZZ_{>0}
 \quad\text{for every pair with }\mu_i+\mu_j<1.
 \label{eq:INT}
\end{equation}
Under this condition the completion by finite-distance collisions has
orbifold universal cover $\mathbb{CH}^{n-3}$, and the monodromy is a
lattice. Cusps remain at infinite distance. For the half-integral
variant, fix a block of equal-weight marks and permit every permutation
within that block. Interchanging two of its marks sends $\delta$ to
$-\delta$; the quotient coordinate is $\delta^2$, so the transverse
angle is halved.
Mostow's $\Sigma$INT condition is therefore
\[
 (1-\mu_i-\mu_j)^{-1}\in
 \begin{cases}
 \tfrac12\ZZ_{>0},&i,j\text{ both in the specified block},\\
 \ZZ_{>0},&\text{otherwise},
 \end{cases}
 \qquad \mu_i+\mu_j<1.
\]
Mostow's theorem gives an orbifold uniformization for that permutation
quotient~\cite[Sections 1--3]{MostowHalfIntegral}; the labeled completion
need not itself be an orbifold when a reciprocal is a nonintegral
half-integer. Thurston gives the corresponding criterion when all
equal-weight permutations are allowed~\cite[Theorem 0.2]{ThurstonShapes}.

When $\mu_i+\mu_j=1$, the collision is a cusp rather than a finite
cone-angle divisor; logarithmic periods and parabolic monodromy enter.
For general real weights satisfying~\eqref{eq:flatweights}, the local
complex hyperbolic geometry and its finite volume still exist
\cite[Theorem 0.2]{ThurstonShapes}. One
must not infer a lattice quotient for arbitrary weights. Conversely,
INT and $\Sigma$INT are the classical sufficient uniformization
conditions, not a claim to classify every exceptional discrete
hypergeometric monodromy group. Our volumes retain all labels; a
quotient by allowed equal-weight permutations changes the quotient
and the volume factors.

\subsection{Four equal deficits and elliptic periods}

For $n=4$ and $\mu_i=1/2$, put the marks at $0,1,z,\infty$. Then
\begin{equation}
 \eta=\frac{\dd w}{\sqrt{w(w-1)(w-z)}},\qquad
 E_z:\ y^2=w(w-1)(w-z).
 \label{eq:ellipticDM}
\end{equation}
The Euclidean cone sphere is the flat pillowcase quotient of this
elliptic curve. The ratio of two periods gives a coordinate
$\tau\in\HH$. For $0<z<1$, one choice is
\begin{equation}
 \tau(z)=i\frac{K(1-z)}{K(z)},\qquad
 K(z)=\int_0^{\pi/2}\frac{\dd u}{\sqrt{1-z\sin^2u}}.
 \label{eq:ellipticperiodratio}
\end{equation}
Here $z$ is the elliptic parameter, not its square root. The inverse
map is the modular lambda function. Projective monodromy is
$\mathrm P\Gamma(2)=\Gamma(2)/\{\pm I\}$, and the labeled moduli
space is the thrice-punctured
sphere with curvature-$-1$ area $2\pi$. The action of $S_4$ on the
cross-ratio has the Klein four-group as its kernel. Its effective
quotient therefore has degree $6$, and is the modular orbifold of
area $\pi/3$; dividing the labeled area by $24$ would instead retain
the ineffective stabilizer
\cite[Example 2.1 and Section 5]{Kojima}. This is the flat
endpoint of the equal-angle torus--sphere correspondence in
Section~\ref{sec:torussphere}.

